\section*{Capítulo \ref{ch:b3:modules} --- Módulos sobre un dominio de
ideales principales}

\begin{solution}{exo:b3:modules:1}
(a) Si $\Q = \Z q_1 + \dots + \Z q_k$, sea $d$ un denominador común de
los $q_i$: toda combinación está en $\frac1d\Z$, pero
$\frac1{2d} \notin \frac1d\Z$. Contradicción.

(b) Sin \omterm{def:b3:modules:torsion}{torsión}: $nq = 0$ con
$n \neq 0$ obliga a $q = 0$ en $\Q$. No
\omterm{def:b3:modules:free}{libre}: dos racionales no nulos
cualesquiera $\frac ab, \frac cd$ cumplen la relación no trivial
$(bc)\frac ab - (ad)\frac cd =
0$, luego una base tiene a lo sumo un elemento; y $\Q \cong \Z$ haría
$\Q = \Z q$ cíclico, pero $\frac q2 \notin \Z q$. (Además $\Q \neq
0$.)

(c) El~\cref{thm:b3:modules:structure} supone la \emph{generación
finita}, que (a) niega: no hay contradicción --- más bien, $\Q$ muestra
que la hipótesis es necesaria en el enunciado
«sin \omterm{def:b3:modules:torsion}{torsión} $\Rightarrow$
\omterm{def:b3:modules:free}{libre}».
\end{solution}

\begin{solution}{exo:b3:modules:2}
$360 = 2^3\cdot3^2\cdot5$. Particiones: de $3$: $(3), (2,1),
(1,1,1)$; de $2$: $(2), (1,1)$; de $1$: $(1)$. Por tanto, $3 \times 2
\times 1 = 6$ grupos. \omterm{thm:b3:modules:structure}{Divisores
elementales} $\to$ \omterm{thm:b3:modules:smith}{factores invariantes}:
\begin{center}
\begin{tabular}{l l}
$\Z/8 \times \Z/9 \times \Z/5$ & $\Z/360$\\
$\Z/8 \times \Z/3 \times \Z/3 \times \Z/5$ & $\Z/3 \times
\Z/120$\\
$\Z/4 \times \Z/2 \times \Z/9 \times \Z/5$ & $\Z/2 \times
\Z/180$\\
$\Z/4 \times \Z/2 \times \Z/3 \times \Z/3 \times \Z/5$ & $\Z/6
\times \Z/60$\\
$(\Z/2)^3 \times \Z/9 \times \Z/5$ & $\Z/2 \times \Z/2 \times
\Z/90$\\
$(\Z/2)^3 \times \Z/3 \times \Z/3 \times \Z/5$ & $\Z/2 \times
\Z/6 \times \Z/30$
\end{tabular}
\end{center}
(Para pasar a \omterm{thm:b3:modules:smith}{factores invariantes}: el
mayor $d_s$ reúne la mayor potencia de cada primo, y así sucesivamente
hacia abajo.) De orden $p^5$: tantos como particiones de $5$, a saber,
$7$.
\end{solution}

\begin{solution}{exo:b3:modules:3}
$B$: $D_1 = \gcd(2,4,6,8) = 2$; $D_2 = \abs{\det B} = \abs{16 -
24} = 8$. \omterm{thm:b3:modules:smith}{Factores invariantes} $d_1 = 2$,
$d_2 = 8/2 = 4$: forma de Smith $\operatorname{diag}(2, 4)$, y
$\Z^2/B\Z^2 \cong \Z/2\Z
\times \Z/4\Z$.

$C$: diagonal, pero \emph{no} de Smith ($2 \nmid 3$). $D_1 =
\gcd(2,3,12) = 1$; $D_2 = \gcd(2\cdot3,\, 2\cdot12,\, 3\cdot12) =
\gcd(6, 24, 36) = 6$; $D_3 = 72$. Luego $d = (1, 6, 12)$ y
$\Z^3/C\Z^3 \cong \Z/6\Z\times\Z/12\Z$ --- coherentemente con el teorema
chino del resto: $\Z/2\times\Z/3\times\Z/12 \cong \Z/6\times\Z/12$.
\end{solution}

\begin{solution}{exo:b3:modules:4}
Escribamos $B = Q\,\operatorname{diag}(d_1, \dots, d_n)\,P$ con $Q, P
\in GL_n(\Z)$ (el~\cref{thm:b3:modules:smith}; ningún $d_i$ es nulo,
pues $\det B \neq 0$). Entonces $\Z^n/B\Z^n \cong \Z^n/
\operatorname{diag}(d)\Z^n = \prod_i \Z/d_i\Z$ (el isomorfismo compuesto
$x \mapsto Q^{-1}x$ de $\Z^n$ lleva $B\Z^n$ sobre
$\operatorname{diag}(d)P\Z^n = \operatorname{diag}(d)\Z^n$). Su cardinal
es $\prod \abs{d_i} = \abs{\det
\operatorname{diag}(d)} = \abs{\det B}$, ya que $\det Q, \det P = \pm
1$. Para $L = \Z(2,0) + \Z(1,3)$: $B = \bigl(\begin{smallmatrix}2
& 1\\ 0 & 3\end{smallmatrix}\bigr)$, $D_1 = 1$, $D_2 = 6$: $\Z^2/L
\cong \Z/6\Z$, de cardinal $\abs{\det B} = 6$.
\end{solution}

\begin{solution}{exo:b3:modules:5}
(a) Submódulo: hecho en la~\cref{def:b3:modules:torsion}. Si $a(x +
T(M)) = 0$ en $M/T(M)$ con $a \neq 0$, entonces $ax \in T(M)$: $bax
= 0$ para cierto $b \neq 0$, y $ba \neq 0$ (dominio), luego $x \in
T(M)$: la clase es nula. Así, $M/T(M)$ es
sin \omterm{def:b3:modules:torsion}{torsión}.

(b) $(X, Y) \subseteq K[X,Y]$ es sin
\omterm{def:b3:modules:torsion}{torsión} (submódulo del dominio $K[X,Y]$ actuando sobre sí mismo).
Supongamos que fuera \omterm{def:b3:modules:free}{libre}; dos elementos
cualesquiera $P, Q$ cumplen $Q\cdot P - P \cdot Q = 0$, una relación no
trivial cuando $P \ne Q$ son no nulos, luego una base tiene un solo
elemento: $(X, Y) = (P)$ sería principal --- en contra
del~\cref{exo:b3:rings:6}(a). Es sin
\omterm{def:b3:modules:torsion}{torsión} y \omterm{def:b3:modules:free}{finitamente generado} ($X, Y$ lo
generan), pero no \omterm{def:b3:modules:free}{libre}: sobre el
no \omterm{def:b3:rings:pidufd}{DIP} $K[X,Y]$, el teorema de estructura
falla.
\end{solution}

\begin{solution}{exo:b3:modules:6}
(a) Si $\Z = 2\Z \oplus C$, la proyección $\Z \to \Z/2\Z$ se restringe a
un isomorfismo $C \cong \Z/2\Z$: $C$ sería un subgrupo de $\Z$ cuyo
elemento no nulo $x$ cumpliría $2x \in C
\cap 2\Z = 0$. Pero $\Z$ es sin
\omterm{def:b3:modules:torsion}{torsión}: $C = 0$, lo que obliga a $\Z =
2\Z$ --- falso.

(b) $A^n/M$ es \omterm{def:b3:modules:free}{finitamente generado} y
sin \omterm{def:b3:modules:torsion}{torsión}, luego
\omterm{def:b3:modules:free}{libre}
(el~\cref{thm:b3:modules:structure}): $A^n/M \cong A^r$ con base
$f_1, \dots, f_r$. Elíjanse preimágenes $y_i \in A^n$ de los $f_i$ y
póngase $F = Ay_1 + \dots + Ay_r$. Todo $x \in A^n$ cumple $\pi(x)
= \sum a_if_i$, luego $x - \sum a_iy_i \in M$: $A^n = M + F$. Si
$\sum a_iy_i \in M$, aplicando $\pi$ se obtiene $\sum a_if_i = 0$ y, por
tanto, todos los $a_i = 0$ (es una base): $M \cap F = 0$. Así
$A^n = M \oplus
F$.
\end{solution}

\begin{solution}{exo:b3:modules:7}
La reducción del~\cref{exo:b3:modules:3} era efectiva: con
\[
L = \begin{pmatrix} 1 & 0\\ -3 & 1\end{pmatrix},
\qquad
R = \begin{pmatrix} 1 & 2\\ 0 & -1 \end{pmatrix},
\qquad
LBR = \begin{pmatrix} 2 & 0\\ 0 & 4\end{pmatrix}
\]
(operación de fila $R_2 \leftarrow R_2 - 3R_1$, operaciones de columna
$C_2 \leftarrow C_2 - 2C_1$ y después $C_2 \leftarrow -C_2$). Poniendo
$y = R^{-1}x$ (una biyección de $(\Z/20\Z)^2$, por ser $R$ invertible
sobre $\Z$), el sistema $Bx \equiv b \pmod{20}$ equivale a
\[
2y_1 \equiv b_1, \qquad 4y_2 \equiv -3b_1 + b_2 \pmod{20}.
\]
La congruencia $ky \equiv c \pmod{20}$ es
\omterm{def:b3:groups:derived}{resoluble} si y solo si
$\gcd(k, 20) \mid c$: hay soluciones si y solo si $2 \mid b_1$ y $4
\mid b_2 - 3b_1$, es decir, \emph{$b_1$ par y $b_2 \equiv 3b_1
\pmod 4$}. Cuando es \omterm{def:b3:groups:derived}{resoluble} hay
$2 \times 4 = 8$ soluciones \omterm{def:b3:modules:module}{módulo} $20$.
\end{solution}

\begin{solution}{exo:b3:modules:8}
(a) Una $u$ nilpotente cumple $\mu_u = X^k$; los
\omterm{thm:b3:modules:structure}{divisores elementales} son
$X^{k_1} \geq \dots$, un bloque de Jordan $J_{k_i}(0)$ por cada parte de
una partición de $4$:
\begin{center}
\begin{tabular}{l l l}
partición & \omterm{thm:b3:modules:jordan}{forma de Jordan} & \omterm{thm:b3:modules:smith}{factores invariantes}\\
$(4)$ & $J_4$ & $X^4$\\
$(3,1)$ & $J_3 \oplus J_1$ & $X,\ X^3$\\
$(2,2)$ & $J_2 \oplus J_2$ & $X^2,\ X^2$\\
$(2,1,1)$ & $J_2 \oplus J_1 \oplus J_1$ & $X,\ X,\ X^2$\\
$(1,1,1,1)$ & $0$ & $X, X, X, X$
\end{tabular}
\end{center}
(b) Tómense $u = J_2\oplus J_2$ y $v = J_2 \oplus J_1 \oplus J_1$: ambas
tienen $\chi = X^4$, $\mu = X^2$, pero
\omterm{thm:b3:modules:smith}{factores invariantes} distintos --- no son
semejantes (el~\cref{thm:b3:modules:frobenius}); también se pueden
comparar los rangos: $\operatorname{rk} u = 2 \neq 1 =
\operatorname{rk} v$. Para $n \leq 3$: $\chi$ y $\mu$ determinan, para
cada valor propio $\lambda$ (sobre un cuerpo de descomposición), el
tamaño total $m_\lambda \leq 3$ de los bloques \omterm{def:b3:rings:divisibility}{asociados} a $\lambda$ y
el mayor bloque $r_\lambda$; y una partición de $m \leq 3$ queda
determinada por su parte mayor ($m = 3, r = 2$ obliga a $(2,1)$, etc.).
Así, los \omterm{thm:b3:modules:structure}{divisores elementales}
coinciden, y el~\cref{cor:b3:modules:descent} hace descender la
semejanza al cuerpo base.
\end{solution}

\begin{solution}{exo:b3:modules:9}
(i)$\Rightarrow$(ii): si $V = K[u]x$, entonces $V \cong
K[X]/\operatorname{Ann}(x)$ y $\operatorname{Ann}(x) = (\mu_u)$ (un
polinomio aniquila $x$ si y solo si aniquila todo $V = K[u]x$, pues
$P(u)Q(u)x = Q(u)P(u)x$). Luego $n = \dim V = \deg \mu_u$; y como
$\mu_u \mid \chi_u$ y $\deg\chi_u = n$, ser mónicos da $\mu_u =
\chi_u$.

(ii)$\Rightarrow$(iii): $\deg\chi_u = \sum_i \deg P_i$ y $\mu_u = P_s$
(el~\cref{thm:b3:modules:frobenius}); la igualdad de grados obliga a
$s = 1$.

(iii)$\Rightarrow$(i): $V \cong K[X]/(P_1)$ es cíclico, generado por la
preimagen de $\bar 1$.

Una \omterm{def:b3:modules:companion}{matriz compañera} es el caso
$V = K[X]/(P)$ mismo: $x = \bar
1$, es decir, $e_1$, es cíclico. Para una matriz diagonal
$\operatorname{diag}(\lambda_1, \dots, \lambda_n)$: $\chi = \prod
(X - \lambda_i)$, $\mu = \prod_{\lambda \text{ distintos}} (X -
\lambda)$; coinciden si y solo si los $\lambda_i$ son distintos dos a
dos: una matriz diagonal tiene vector cíclico si y solo si sus entradas
diagonales son distintas dos a dos (entonces sirve $x = (1, \dots, 1)$:
Vandermonde).
\end{solution}

\begin{solution}{exo:b3:modules:10}
Smith: $M = Q\operatorname{diag}(d_1,\dots,d_n)P$;
el~\cref{exo:b3:modules:4} da $[\Z^n : M\Z^n] = \prod\abs{d_i} =
\abs{\det M}$. Si $\det M = \pm 1$: la fórmula de la matriz adjunta
$M^{-1} = (\det M)^{-1}\,{}^{t}\!\operatorname{com}(M)$ tiene entradas
enteras, luego $M \in GL_n(\Z)$; recíprocamente, $MM^{-1} = I$ da
$\det M \cdot \det M^{-1} = 1$ en $\Z$, luego $\det M = \pm1$.

Subgrupos de índice $m$ en $\Z^2$: un tal subgrupo $L$ tiene rango $2$
(índice finito) y una única base en \emph{forma normal de Hermite}
$\bigl(\begin{smallmatrix} a & b\\ 0 &
d\end{smallmatrix}\bigr)$: $d$ queda caracterizado por $L \cap (\{0\}
\times \Z) = \{0\} \times d\Z$, y $a$ por $\pi_1(L) = a\Z$ (primeras
coordenadas), y $b$ es entonces único \omterm{def:b3:modules:module}{módulo} $d$; normalícense $a, d
> 0$ y $0 \leq b < d$. El índice es $ad = m$. Recuento: para cada
divisor $d \mid m$ ($a = m/d$) hay $d$ elecciones de $b$: en total
$\sum_{d \mid m} d = \sigma_1(m)$.
\end{solution}

\begin{solution}{exo:b3:modules:11}
(a) Por el teorema de estructura, $G \cong \prod_i\Z/d_i\Z$, y $x^k = e$
se desacopla coordenada a coordenada. En $\Z/d\Z$: $kx \equiv 0
\pmod d$ tiene exactamente $\gcd(k, d)$ soluciones ($x$ ha de ser
múltiplo de $d/\gcd(k,d)$, y hay $\gcd(k, d)$ de ellos). Multiplíquese
sobre los factores.

(b) Si $G = \Z/d_s\Z$ es cíclico ($s = 1$), el recuento es
$\gcd(k, d_s) \leq k$. Si $s \geq 2$: tómese $k = d_1$; el recuento es
$\prod_i\gcd(d_1, d_i) = d_1^{\,s} > d_1$ (cada $\gcd$ vale $d_1$ por la
cadena de divisibilidad): la ecuación $x^{d_1} = e$ tiene más de $d_1$
soluciones.

(c) En un cuerpo, $X^k - 1$ tiene a lo sumo $k$ raíces
(el~\cref{ch:b3:rings}: un polinomio no nulo de grado $k$ sobre un
dominio), de modo que todo subgrupo finito $G \leq K^\times$ cumple el
criterio de (b): $G$ es cíclico --- la demostración estructural, de una
línea, de la ciclicidad de $\mathbb F_q^\times$, complementaria de la
demostración por recuento de
\cref{ch:b3:galois}.
\end{solution}

\begin{solution}{exo:b3:modules:12}
(a) Si $MN = I$ con $N$ enteras: $\det M\det N = 1$ con ambos enteros,
luego $\det M = \pm1$. Recíprocamente, si $\det M =
\pm1$, la fórmula de los cofactores $M^{-1} =
\frac1{\det M}\,{}^t\!\operatorname{com}(M)$ tiene entradas enteras.

(b) Multiplicar por la izquierda por $E^{-k}$ resta $k$ veces la fila
$2$ a la fila $1$; por $F^{-k}$, $k$ veces la fila $1$ a la fila $2$.
Dada $M = \bigl(\begin{smallmatrix}a & *\\ c &
*\end{smallmatrix}\bigr) \in SL_2(\Z)$, la primera columna $(a, c)$ es
un vector unimodular ($\gcd(a, c) = 1$: divide a $\det M = 1$).
Ejecútese Euclides sobre $(a, c)$ con estas operaciones de fila: tras
finitos pasos la columna se convierte en $(\pm1, 0)$. La matriz es ahora
$\pm
\bigl(\begin{smallmatrix}1 & b\\ 0 & 1\end{smallmatrix}
\bigr) = \pm E^{b}$ (el determinante sigue siendo $1$). Queda escribir
$-I$ en los generadores: $(EF^{-1}E)^2 =
\bigl(\begin{smallmatrix}0 & 1\\ -1 &
0\end{smallmatrix}\bigr)^2 = -I$ (compruébese el cuadrado de la matriz
de rotación). Deshaciendo el proceso, $M$ es una palabra en $E^{\pm1},
F^{\pm1}$.

(c) El algoritmo de Smith (el~\cref{met:b3:modules:smith}) usa
exactamente esas operaciones de filas y columnas (más intercambios y
cambios de signo, que a su vez son productos de operaciones elementales
salvo el signo del determinante): sobre $\Z$, toda matriz es $U\,D\,V$
con $U, V$ productos de matrices elementales y $D$ la forma de Smith ---
(b) es el caso $2\times2$ de determinante $1$ del hecho general de que
las matrices de tipo $E$ generan
$SL_n(\Z)$.
\end{solution}

\begin{solution}{pb:b3:modules:1}
\textbf{1.} $\varphi$ es aditiva y $K[X]$-lineal: $\varphi(X
\cdot (Q_i)_i) = \sum_i (XQ_i)(M)e_i = M\sum_i Q_i(M)e_i = X
\cdot \varphi\bigl((Q_i)_i\bigr)$, siendo $X \cdot v = Mv$ la estructura
de \omterm{def:b3:modules:module}{módulo} de $V_M$. Es sobreyectiva: los
vectores constantes dan todo $K^n$. La columna $j$ de $XI - M$ es
$Xe_j - \sum_i
m_{ij}e_i$, cuya imagen es $Me_j - \sum_i m_{ij}e_i = 0$.

\textbf{2.} \omterm{def:b3:modules:module}{Módulo} las columnas, $Xe_j \equiv \sum_i m_{ij}e_i$: todo
vector de polinomios se reduce, por inducción sobre el grado máximo, a
un vector \emph{constante} $c \in K^n$. Si el vector de partida está en
$\ker\varphi$, entonces $\varphi(c) = c = 0$ (sobre las constantes,
$\varphi$ es la identificación $K^n = V_M$), de modo que el vector está
en el subespacio generado por las columnas: $\ker\varphi = (XI_n -
M)K[X]^n$. Por tanto $V_M \cong K[X]^n/(XI - M)K[X]^n$, y Smith sobre
$K[X]$ (todos los \omterm{thm:b3:modules:smith}{factores invariantes}
son no nulos, pues su producto es $\det(XI - M) = \chi_M$) da
$V_M \cong \bigoplus_i
K[X]/(f_i)$: los $f_i$ no constantes son los
\omterm{thm:b3:modules:frobenius}{invariantes de semejanza}, calculables
como cocientes de máximos comunes divisores de menores
(\cref{thm:b3:modules:smith}).

\textbf{3.} $\lambda I_n$: $XI - \lambda I$ ya está en forma de Smith:
invariantes $(X - \lambda, \dots, X - \lambda)$, $n$ de ellos. Entradas
diagonales distintas: $V \cong \bigoplus_i K[X]/(X -
\lambda_i)$ con \omterm{def:b3:modules:module}{módulos} \omterm{thm:b3:rings:crt}{comaximales} dos a
dos, de modo que el teorema chino del resto lo comprime en el único
cíclico $K[X]/\bigl(\prod_i(X - \lambda_i)\bigr)$: un solo invariante,
$\chi$. $J_n(0)$: $\mu = X^n = \chi$ obliga a un único invariante $X^n$.
$\operatorname{diag}(J_2(0), J_1(0))$:
\omterm{thm:b3:modules:structure}{divisores elementales} $X^2, X$;
invariantes $P_1 = X \mid P_2 =
X^2$.

\textbf{4.} $P \mapsto P(u)$ aplica $K[X]$ sobre $K[u]$, con núcleo
$(\mu_u)$ por definición del polinomio mínimo: $K[u]
\cong K[X]/(\mu_u)$, de dimensión $\deg\mu_u = \deg P_s = n_s$.

\textbf{5.} La $K[X]$-linealidad obliga a $f(\bar Q) = f(Q\cdot\bar 1)
= Q\,f(\bar 1)$. La clase $\bar 1$ cumple $P \bar 1 = 0$, luego
$P f(\bar 1) = 0$ es necesario. Recíprocamente, si $P\bar c = 0$,
entonces $f(\bar Q) = Q\bar c$ está bien definida ($Q \equiv Q' \bmod P
\Rightarrow (Q - Q')\bar c \in$ múltiplos de $P\bar c = 0$) y es
$K[X]$-lineal.

\textbf{6.} Sean $g = \gcd(P, Q)$, $P = gP'$, $Q = gQ'$ con
$\gcd(P', Q') = 1$. En $K[X]/(Q)$: $P\bar c = 0 \iff Q \mid Pc
\iff Q' \mid P'c \iff Q' \mid c$ (Euclides, $\gcd(P', Q') = 1$). Así
pues, los $\bar c$ admisibles forman el submódulo generado por
$\bar{Q'}$, cuyo anulador es $\{R : Q \mid RQ'\} = (g)$: ese submódulo
es $\cong K[X]/(g)$. Con la pregunta 5,
$\operatorname{Hom}(K[X]/(P), K[X]/(Q)) \cong K[X]/(\gcd(P,Q))$,
de dimensión $\deg\gcd(P,Q)$.

\textbf{7.} $v$ conmuta con $u$ si y solo si $v$ conmuta con todo
$P(u)$, si y solo si $v$ es $K[X]$-lineal: $\mathcal C(u) =
\operatorname{End}_{K[X]}(V)$. Escribiendo los morfismos de $V =
\bigoplus_j V_j$ como matrices $(f_{ij})$, $f_{ij} \in
\operatorname{Hom}(V_j, V_i)$ (compóngase con las inyecciones y las
proyecciones), la pregunta 6 da
\[
\dim \mathcal C(u) = \sum_{i,j} \deg\gcd(P_i, P_j)
= \sum_{i,j} n_{\min(i,j)}
= \sum_{k=1}^{s} \bigl(2(s - k) + 1\bigr)\,n_k ,
\]
usando la cadena de divisibilidad ($\gcd(P_i, P_j) =
P_{\min(i,j)}$) y, para el último paso, que $\min(i,j) = k$ ocurre
exactamente para $2(s - k) + 1$ pares $(i, j)$.

\textbf{8.} Como $2(s-k)+1 \geq 1$, se tiene $\dim\mathcal C(u) \geq
\sum_k n_k = n$, con igualdad si y solo si $s = 1$, es decir, si y solo
si $u$ es cíclico (el~\cref{exo:b3:modules:9}). Para
$u = \lambda\,\mathrm{id}$: $s = n$, todos $n_k = 1$:
$\dim = \sum_{k=1}^n (2(n-k)+1) = n^2$ --- correcto, pues
$\mathcal C(\lambda\,\mathrm{id}) = \mathcal
L(V)$.

\textbf{9.} Fórmula: invariantes $(X, X^2)$, luego $s = 2$, $n_1 =
1$, $n_2 = 2$: $\dim = 3\cdot1 + 1\cdot2 = 5$. Directamente: en la base
$(e_1, e_2, e_3)$ con $u e_2 = e_1$, $ue_1 = ue_3 = 0$, escribiendo
$Au = uA$ para $A = (a_{ij})$ se obtienen las condiciones
$a_{21} = a_{23} = a_{31} = 0$ y $a_{11} = a_{22}$: cinco parámetros
\omterm{def:b3:modules:free}{libres}
$a_{11}{=}a_{22}, a_{12}, a_{13}, a_{32}, a_{33}$.

\textbf{10.} Un polinomio $P(u)$ conmuta con todo lo que conmute con $u$
(es una suma de potencias de $u$): $K[u]
\subseteq \mathcal C(\mathcal C(u))$. Y $u \in \mathcal C(u)$, de modo
que toda $w \in \mathcal C(\mathcal C(u))$ conmuta con $u$.

\textbf{11.} Sean $V = K[u]x$ y $v \in \mathcal C(u)$. Escribamos
$v(x) = P(u)x$ (ciclicidad). Para $y = Q(u)x$ arbitrario: $v(y) =
vQ(u)x = Q(u)v(x) = Q(u)P(u)x = P(u)y$. Luego $v = P(u)$: $\mathcal
C(u) = K[u]$. Entonces $\mathcal C(\mathcal C(u)) = \mathcal
C(K[u]) = \mathcal C(u) = K[u]$ (conmutar con todo $K[u]$ es lo mismo
que conmutar con $u$). El teorema vale en el caso cíclico.

\textbf{12.} $\pi_i$ es $K[X]$-lineal (la descomposición es una suma
directa de submódulos), luego $\pi_i \in \mathcal C(u)$, y $w$ conmuta
con ella: $w(V_i) = w\pi_i(V) = \pi_i w(V) \subseteq
V_i$. La restricción $w_i = w\restriction_{V_i}$ conmuta con el
$u_i = u\restriction_{V_i}$ cíclico (pregunta 10), y $w_i
\in \mathcal C(u_i) = K[u_i]$ (pregunta 11): $w_i = Q_i(u_i) =
Q_i(u)\restriction_{V_i}$ para cierto $Q_i \in K[X]$.

\textbf{13.} Buena definición de $\eta_{ij}(R(u)x_j) = R(u)x_i$: si
$R(u)x_j = 0$, entonces $P_j \mid R$, y $P_i \mid P_j$ da $P_i \mid R$,
luego $R(u)x_i = 0$ ($\operatorname{Ann}(x_i) =
(P_i)$). $\eta_{ij}$ es entonces $K[X]$-lineal por construcción, y
$\tilde\eta_{ij} = \eta_{ij}\pi_j$ es una composición de
$K[X]$-morfismos $V \to V$: $\tilde\eta_{ij} \in \mathcal
C(u)$.

\textbf{14.} Evaluemos $w\tilde\eta_{ij} = \tilde\eta_{ij}w$ en $x_j$:
el miembro de la izquierda es $w(x_i) = Q_i(u)x_i$; el de la derecha,
$\eta_{ij}\bigl(Q_j(u)x_j\bigr) = Q_j(u)x_i$. Por tanto $(Q_i -
Q_j)(u)\,x_i = 0$: $P_i \mid Q_i - Q_j$ para todo $i \leq j$. En
particular, con $Q = Q_s$: $Q \equiv Q_i \pmod{P_i}$, luego $Q(u)$ y
$Q_i(u)$ coinciden sobre $V_i$ (al que $P_i(u)$ aniquila). Así pues, $w
= Q(u)$ sobre cada $V_i$ y, por tanto, sobre $V$:
$\mathcal C(\mathcal C(u))
\subseteq K[u]$ y, con la pregunta 10, $\mathcal C(\mathcal
C(u)) = K[u]$.

\textbf{15.} La primera afirmación son las preguntas 10--14 (o, para $u$
cíclico, la pregunta 11 sola). Para $u = \mathrm{id}$, $n \geq
2$: $\mathcal C(u) = \mathcal L(V)$ tiene dimensión $n^2$, mientras que
$\dim K[u] = \deg\mu_u = 1$. Así que $\mathcal C(u) = K[u]$ falla
estrepitosamente; y sin embargo el teorema del biconmutante se cumple
($\mathcal C(\mathcal L(V)) = K\,\mathrm{id} = K[u]$: el
\omterm{ex:b3:groups:actions}{centro} del álgebra de matrices son los
escalares). La ciclicidad es lo que hace que el conmutante \emph{simple}
ya sea polinómico; el \emph{doble} conmutante lo es siempre.

\textbf{16.} Si $P(XI - M)Q = S$ es una reducción de Smith ($P, Q$
invertibles sobre $K[X]$), trasponiendo se obtiene ${}^tQ\,(XI -
{}^tM)\,{}^tP = {}^tS = S$: la misma forma de Smith, luego $XI - M$ y
$XI - {}^tM$ tienen los mismos \omterm{thm:b3:modules:smith}{factores
invariantes}, es decir, $M$ y ${}^tM$ tienen los mismos
\omterm{thm:b3:modules:frobenius}{invariantes de semejanza}
(el~\cref{cor:b3:modules:descent}): son semejantes.

\textbf{17.} Los \omterm{thm:b3:modules:frobenius}{invariantes de
semejanza} de $M$ sobre $L$ son los
\omterm{thm:b3:modules:smith}{factores invariantes} de $XI - M$ en
$L[X]$. Una reducción de Smith de $XI - M$ sobre $K[X]$ --- con $P, Q$
invertibles sobre $K[X]$ y forma diagonal con la cadena de divisibilidad
--- es \emph{también} una reducción de Smith válida sobre $L[X]$ ($P, Q$
siguen siendo invertibles: sus determinantes son constantes no nulas), y
los \omterm{thm:b3:modules:smith}{factores invariantes} mónicos son
únicos: los \omterm{thm:b3:modules:smith}{factores invariantes}
calculados sobre $K$ y sobre $L$ coinciden. Así, $M \sim_L N$ si y solo
si tienen los mismos \omterm{thm:b3:modules:smith}{factores
invariantes}, si y solo si $M \sim_K N$. En particular, dos matrices
reales conjugadas en $\C$ lo son en $\R$ --- un enunciado que suele
demostrarse por vía analítica (especializando una $P + \iu
Q$ invertible) y aquí sale por vía estructural.

\textbf{18.} Descompóngase $u = \bigoplus J_{m_t}(0)$ en bloques de
Jordan nilpotentes. En un bloque de tamaño $m$,
$\operatorname{rk}J_m^k = \max(m - k, 0)$, luego
$\operatorname{rk}J_m^{k-1} - \operatorname{rk}J_m^k = 1$ si $m
\geq k$ y $0$ en caso contrario. Sumando sobre los bloques:
$\operatorname{rk}u^{k-1} - \operatorname{rk}u^k =
\#\{t : m_t \geq k\}$. La sucesión de rangos determina, pues, el
multiconjunto $(m_t)$ --- una partición de $n$ --- y recíprocamente toda
partición se realiza: las clases de nilpotentes $\leftrightarrow$
particiones de $n$, sobre cualquier cuerpo. $M_5$: $p(5) = 7$ clases
($5$; $4{+}1$; $3{+}2$; $3{+}1{+}1$; $2{+}2{+}1$;
$2{+}1{+}1{+}1$; $1^5$).

\textbf{19.} (a) Sobre $\Q$ (o sobre cualquier cuerpo de característica
$0$), $D^n = 0$, $D^{n-1}(X^{n-1}) = (n-1)!\, \neq 0$: $D$ es nilpotente
de índice $n$ sobre un espacio de dimensión $n$, luego es cíclico con
único invariante $X^n$ ($x = X^{n-1}$ lo genera: sus derivadas sucesivas
generan el espacio). (b) Sobre $\mathbb F_p$ con $p
< n$: $D^p = 0$, porque la derivada $p$-ésima de todo monomio $X^m$
lleva el factor $m(m-1)\cdots(m-p+1)$, producto de $p$ enteros
consecutivos y, por tanto, $\equiv 0 \bmod p$. Escribamos $n = ap + r$,
$0 \leq r < p$. Entonces $\ker D^k$ está generado por los monomios $X^m$
con $D^kX^m = 0$; contando los exponentes $m < n$ según su resto \omterm{def:b3:modules:module}{módulo}
$p$: $\dim\ker D^k = ak + \min(r,
k)$ para $0 \leq k \leq p$, luego $\operatorname{rk}D^{k-1} -
\operatorname{rk}D^k = \dim\ker D^k - \dim\ker D^{k-1} = a +
\mathbf 1_{k \leq r}$. Por la pregunta 18, la partición tiene $a$
bloques de tamaño exactamente $p$ y (si $r > 0$) un bloque de tamaño
$r$: \omterm{thm:b3:modules:frobenius}{invariantes de semejanza}
$X^r \mid X^p \mid \dots \mid X^p$. La característica cambia la forma
canónica del operador más familiar de las matemáticas.

\textbf{20.} $M \in GL_2$ tiene $s \in \{1, 2\}$
\omterm{thm:b3:modules:smith}{factores invariantes}. Si $s = 2$:
$P_1 = P_2 = X - a$ ($a \neq 0$: invertibilidad), es decir, $M = aI$,
central. Si $s = 1$: $M$ es cíclica, con polinomio característico $=$
igual al polinomio mínimo $\chi$ de grado $2$, y las clases corresponden
a los posibles $\chi$ con $\chi(0) \neq 0$: $\chi$ descompuesto con
raíces distintas $a \neq
b$ (la compañera $\sim$ es diagonalizable); $\chi = (X - a)^2$ (la
compañera, no semisimple); $\chi$
\omterm{def:b3:rings:divisibility}{irreducible}. Exactamente un tipo
cada uno --- los \omterm{thm:b3:modules:smith}{factores invariantes} son
un invariante completo.

\textbf{21.} Centrales: $q - 1$ elecciones de $a$. Valores propios
distintos y descompuestos: pares no ordenados $\{a, b\} \subseteq
\mathbb F_q^\times$, $a \neq b$: $\binom{q-1}2$ clases. Polinomio mínimo
$(X-a)^2$: $q - 1$ clases. Cuadráticas
\omterm{def:b3:rings:divisibility}{irreducibles} con término
independiente no nulo: todas las cuadráticas
\omterm{def:b3:rings:divisibility}{irreducibles} valen (sus raíces son
no nulas), y hay $\frac{q^2 - q}2$ cuadráticas mónicas
\omterm{def:b3:rings:divisibility}{irreducibles} (las $q^2$ cuadráticas
mónicas menos las $\binom q2 + q = \frac{q^2+q}2$ que se descomponen).
En total:
\[
(q - 1) + \frac{(q-1)(q-2)}2 + (q - 1) + \frac{q^2 - q}2
= q^2 - 1 .
\]

\textbf{22.} Si $M$ no es cíclica, $s = 2$ y $M$ es escalar (la
dicotomía de la pregunta 20 vale en $M_2$, sea invertible o no: dos
\omterm{thm:b3:modules:smith}{factores invariantes} de grado $1$ con
$P_1 \mid P_2$ y $P_1 =
P_2$ obligan a $M = aI$). Las escalares son $q$ entre las $q^4$
matrices: probabilidad de ser cíclica $1 - q^{-3}$. En general, el lugar
no cíclico de $M_n$ es donde los menores $(n-1)\times(n-1)$ de $XI - M$
comparten un factor --- una condición algebraica propia ---, de modo que
su proporción es $O(1/q)$-pequeña para $q$ grande: las matrices con
$\chi = \mu$ son la regla, y el argumento de pegado de la Parte III es
el precio que se paga por las excepciones.

\textbf{23.} Un elemento del \omterm{ex:b3:groups:actions}{centro} de
$\mathcal C(u)$ está en $\mathcal C(u)$ y conmuta con todo elemento de
$\mathcal C(u)$, es decir, está en $\mathcal C(\mathcal C(u)) =
K[u]$ (pregunta 14). Recíprocamente, $K[u] \subseteq \mathcal C(u)$, y
todo $P(u)$ conmuta con todo $v \in \mathcal C(u)$ (un tal $v$ conmuta
con $u$ y, por tanto, con cada potencia de $u$): $K[u]$ es central en
$\mathcal C(u)$. Por tanto $Z(\mathcal C(u))
= K[u]$. En consecuencia, $\mathcal C(u)$ es conmutativa si y solo si
$\mathcal C(u) = Z(\mathcal C(u)) = K[u]$; en ese caso
$\dim\mathcal C(u) = \dim K[u] = n_s \leq n$, mientras que la pregunta 8
da $\dim\mathcal C(u) \geq n$: luego $n_s = n$ y $s = 1$, es decir, $u$
es cíclico. Recíprocamente, para $u$ cíclico la pregunta 11 da
$\mathcal C(u) = K[u]$, que es conmutativa. Estructuralmente: un álgebra
de matrices igual a su propio \omterm{ex:b3:groups:actions}{centro} es
exactamente un álgebra de polinomios $K[u]$ de un $u$ cíclico.

\textbf{24.} Cada coeficiente $2s - 2i + 1$ de la fórmula de Frobenius
es impar, luego
\[
\dim\mathcal C(u) = \sum_{i=1}^s(2s - 2i + 1)\,n_i
\equiv \sum_{i=1}^s n_i = n \pmod 2 .
\]
Para $n = 4$, las sucesiones de grados $n_1 \leq \dots
\leq n_s$ posibles de los \omterm{thm:b3:modules:smith}{factores
invariantes}, con suma $4$, son $(4)$, $(1, 3)$, $(2, 2)$, $(1, 1, 2)$,
$(1, 1, 1, 1)$; todas se realizan, por ejemplo con el nilpotente de
$P_i = X^{n_i}$ (la cadena de divisibilidad se cumple automáticamente).
La fórmula da, respectivamente,
\[
1\cdot4 = 4, \quad 3 + 3 = 6, \quad 6 + 2 = 8, \quad
5 + 3 + 2 = 10, \quad 7 + 5 + 3 + 1 = 16 .
\]
Así pues, el conjunto de valores es $\{4, 6, 8, 10, 16\}$: números pares
de la paridad correcta, pero $12$ y $14$ no aparecen nunca --- entre las
sucesiones casi cíclicas y el $n^2$ de la escalar hay un salto.

\textbf{25.} $\abs{GL_2(\mathbb F_3)} = (q^2 - 1)(q^2 - q) =
8 \cdot 6 = 48$. Tipo central: $I$ y $2I$, dos clases de tamaño $1$. En
los otros tres tipos, $M$ es cíclica (pregunta 20), luego su
centralizador en $GL_2$ es el grupo de los elementos invertibles de
$\mathcal C(M) = K[M]$ (pregunta 11), y el tamaño de la clase es
$=48/\abs{K[M]^\times}$ por órbita--estabilizador. Valores propios
distintos y descompuestos: solo el par $\{1, 2\}$, una clase; $K[M]
\cong \mathbb F_3 \times \mathbb F_3$ (teorema chino del resto sobre
$\chi = (X-1)
(X-2)$), unidades $2 \cdot 2 = 4$, tamaño $48/4 = 12$. Polinomio mínimo
$(X - a)^2$, $a \in \{1, 2\}$: dos clases; $K[M] \cong
\mathbb F_3[X]/((X-a)^2)$, unidades $q^2 - q = 6$ (término independiente
de la unidad $\neq 0$ tras centrar), tamaño $48/6 = 8$. $\chi$
\omterm{def:b3:rings:divisibility}{irreducible}: las cuadráticas mónicas
\omterm{def:b3:rings:divisibility}{irreducibles} sobre $\mathbb F_3$ son
$(9 - 3)/2 = 3$, a saber,
\[
X^2 + 1, \qquad X^2 + X + 2, \qquad X^2 + 2X + 2,
\]
(sin raíces en $\mathbb F_3$: compruébese $0, 1, 2$); tres clases,
$K[M] \cong \mathbb F_9$, unidades $q^2 - 1 = 8$, tamaño $48/8 =
6$. \omterm{cor:b3:groups:classeq}{Ecuación de clases}:
$2\cdot1 + 1\cdot12 + 2\cdot8 + 3\cdot6 =
2 + 12 + 16 + 18 = 48$; y $2 + 1 + 2 + 3 = 8 = q^2 - 1$ clases, en
concordancia con la pregunta 21.
\end{solution}
